¥thispagestyle{empty}

¥noindent
¥textbf{¥Large 著者について}

¥vspace*{3mm}

¥hrule width 450pt height 2pt

¥noindent
¥begin{tabular}{lll}
1968 年 & 2 月 & 埼玉県に生まれる¥¥
1990 年 & 3 月 & 東京理科大学 理学部 物理学科 卒業¥¥
1992 年 & 3 月 & 東北大学 大学院 理学研究科 原子核理学専攻 博士課程前期 2 年の課程 終了¥¥
1996 年 & 3 月 & 東北大学 大学院 理学研究科 原子核理学専攻 博士課程後期 3 年の課程 終了¥¥
1996 年 & 12 月 & 郵政省 認可法人 通信・放送機構 国内招へい研究者として、
東北大学 加齢¥¥
& & 医学研究所に勤務¥¥
1997 年 & 4 月 & 青森大学 工学部 情報システム工学科 助手¥¥
現在 & & 同上
¥end{tabular}

¥setlength{¥baselineskip}{6mm}
¥vspace*{0.5cm}

UNIX にはじめて触れたのは、 22 才のとき、仙台のソフトウェア開発会社
(株)コムテック
(¥texttt{http://www}
¥texttt{.comtec.co.jp/})
でのアルバイトのこと。
そこには、 UNIX System V のマニュアルが一揃い完備していて、
とても勉強になった。

しばらくは、 System V 系のシステムとの日々が続くが、
後に SunOS で BSD 系のシステムにも触れることになる。
 SunOS では、 make 三昧の日々が続き、すっかり UNIX 系のシステムに
はまることになる。
現在では、 FreeBSD 以外に SPARCStation で OpenBSD を走らせる日々。

趣味は、本を読むことと、絵を描くこと、本を作ることなど。
絵を描く時間が取れないのが、最近の悩み。
なお、この本は ASCII の日本語 p¥LaTeX を FreeBSD や Cygwin の上で
使用して作成した。
ちなみにある種の本を作るときも、全く同様の環境を使っているらしい。

¥vspace*{1.2cm}

¥noindent
¥textbf{¥Large はじめての BSD rev.}

¥vspace*{3mm}

¥hrule width 450pt height 2pt

¥noindent
¥begin{tabular}{lll}
2002 年 & 9 月 25 日 & 初版
¥end{tabular}

¥begin{description}
¥item[著者:] ¥ruby{小久保}{こくぼ} ¥ruby{温}{あつし}

〒 030-0943 青森市 幸畑 2-3-1 青森大学 工学部 情報システム工学科

¥texttt{TEL: 017-738-2001 (青森大学代表) / FAX: 017-738-2030 (青森大学代表)}

¥texttt{e-mail: kokubo@aomori-u.ac.jp}

¥texttt{web page: http://www.aomori-u.ac.jp/staff/inform/kokubo/}

¥hspace*{2.2cm}¥texttt{http://www.dma.aoba.sendai.jp/¥symbol{"7E}acchan/}
¥item[発行:] 青森大学出版局

¥item[印刷:] 青森コロニー印刷

〒 030-0943 青森市 幸畑 字松元 62-3 ~~~ ¥texttt{TEL: 017-738-2021 / FAX: 017-738-6753}
¥end{description}

¥hrule width 450pt height 2pt
¥vspace*{1mm}

¥noindent
¥textbf{¥LARGE A 1st Step of BSD rev.}~~~by Atsushi Kokubo

¥noindent
¥copyright Atsushi Kokubo 2002, Printed in Japan.

